\slajdid{06-s021}
\begin{frame}[label=06-s021]{Grejov kod – prelazak „u“ i „iz“}
\gusttekst\centering Računski\par\vspace{4mm}
\begin{columns}[c,onlytextwidth]
\column{.40\textwidth}\centering Uočiti\quad $(b\oplus p)\oplus p=b$\par\vspace{3mm}
\renewcommand{\arraystretch}{1.4}\begin{tabular}{ccccc}\toprule
\textcolor{DEtekst}{$b_{n-1}$} & \textcolor{DEtekst}{$b_{n-2}$} & \textcolor{DEtekst}{$\ldots$} & \textcolor{DEtekst}{$b_{1}$} & \textcolor{DEtekst}{$b_{0}$}\\
\textcolor{DEtekst}{$\oplus$} & \textcolor{DEtekst}{$\oplus$} & \textcolor{DEtekst}{$\oplus$} & \textcolor{DEtekst}{$\oplus$} & \textcolor{DEtekst}{$\oplus$}\\
\textcolor{DEplava}{$p_{n-1}$} & \textcolor{DEplava}{$p_{n-2}$} & \textcolor{DEplava}{$\ldots$} & \textcolor{DEplava}{$p_{1}$} & \textcolor{DEplava}{$p_{0}$}\\
\midrule
\textcolor{DEcrvena}{$a_{n-1}$} & \textcolor{DEcrvena}{$a_{n-2}$} & \textcolor{DEcrvena}{$\ldots$} & \textcolor{DEcrvena}{$a_{1}$} & \textcolor{DEcrvena}{$a_{0}$}\\
\textcolor{DEtekst}{$\oplus$} & \textcolor{DEtekst}{$\oplus$} & \textcolor{DEtekst}{$\oplus$} & \textcolor{DEtekst}{$\oplus$} & \textcolor{DEtekst}{$\oplus$}\\
\textcolor{DEplava}{$p_{n-1}$} & \textcolor{DEplava}{$p_{n-2}$} & \textcolor{DEplava}{$\ldots$} & \textcolor{DEplava}{$p_{1}$} & \textcolor{DEplava}{$p_{0}$}\\
\midrule
\textcolor{DEtekst}{$b_{n-1}$} & \textcolor{DEtekst}{$b_{n-2}$} & \textcolor{DEtekst}{$\ldots$} & \textcolor{DEtekst}{$b_{1}$} & \textcolor{DEtekst}{$b_{0}$}\\
\bottomrule\end{tabular}

\column{.57\textwidth}
\begin{align*}
\text{binarni}\quad B&=b_{n-1}b_{n-2}b_{n-3}\ldots b_1\,b_0\\
\text{Grejov}\quad G&=g_{n-1}g_{n-2}g_{n-3}\ldots g_1\,g_0
\end{align*}
\begin{columns}[T,onlytextwidth]\column{.48\textwidth}
\begin{align*}
g_{n-1}&=b_{n-1}\\g_{n-2}&=b_{n-1}\oplus b_{n-2}\\g_{n-3}&=b_{n-2}\oplus b_{n-3}\\&\ldots\\g_1&=b_2\oplus b_1\\g_0&=b_1\oplus b_0
\end{align*}
\column{.49\textwidth}
\begin{align*}
b_{n-1}&=g_{n-1}\\b_{n-2}&=g_{n-1}\oplus g_{n-2}\\b_{n-3}&=g_{n-2}\oplus g_{n-3}\\&\ldots\\b_1&=g_2\oplus g_1\\b_0&=g_1\oplus g_0
\end{align*}\end{columns}\end{columns}
\end{frame}
